\section{Arithmetic-progression restrictions and differential comparison}
\label{ap:section}

An attainable value contains information about nearby source laws, not merely
about a single tensor restriction.  This section makes that information
quantitative.  If a supporting plane loses too little value when the mean
grade changes, one can steer long grade sequences into many disjoint count
windows at a total cost smaller than the entropy of the windows.  Their
labels then give an improved tensor construction, contradicting the supporting
plane.  The steering argument must control every path permitted by its controls:
an estimate that only holds with high probability would not justify counting
all the output tensors.  We first construct the shared count labels and
prove their finite tensor gain.  We then build controls whose cost is
strictly smaller than that gain.

All source laws in this section are the canonical laws in each of the six
physical orientation banks of Definition~\ref{def:strong}.  In particular,
they are not the aggregate physical law obtained by forgetting the banks.

Let
\[
 \mathcal P=\{x\in\R^3:x_0+x_1+x_2=0\}.
\]
We use orthonormal coordinates on this plane when forming covariances,
gradients, and Hessians.  For the six-state alphabet, the grade of a source
letter is its vector in
\(S=\{2e_i,\boldsymbol 1-e_i:0\leq i\leq2\}\).
For a law \(\alpha\), write
\[
 \mu(\alpha)=\E_\alpha g,
 \qquad
 C(\alpha)=\E_\alpha[(g-\mu)(g-\mu)^\top]\big|_{\mathcal P}.
\]
The covariance is positive definite when every letter has positive
probability.  Indeed, the three grades \(2e_i\) alone affinely span the mean
plane.  For a twice continuously differentiable chart \(\alpha(\mu)\)
whose grade mean is \(\mu\), set
\[
 \mathcal L u=C(\alpha(\mu)):D^2u
              =\Tr(C(\alpha(\mu))D^2u).
\]
This is the full covariance contraction; the corresponding probabilistic
generator has an additional factor \(1/2\).

Smooth touching tests are a standard formulation in viscosity-solution
theory \cite{crandalllions1983}; see in particular
\cite[Section~2]{crandallishiilions1992}.  The differential inequality
below is derived directly from finite tensor constructions.

\begin{theorem}[Lower tests for the six-state value]\label{thm:ap}
Fix \(0<\beta\leq3\), and let \(\alpha(\mu)\) be a \(C^2\) chart of
strictly positive six-state laws, parametrized by their means on an open
subset of \(\{\mu:\sum_i\mu_i=2\}\).  Suppose that a \(C^2\) function
\(\varphi\) touches \(F_\beta\circ\alpha\) from below at an interior
point \(\mu_*\).  If \(g_*\) is any supporting supergradient of
\(F_\beta\) at \(\alpha(\mu_*)\), then
\begin{equation}\label{ap:lower-test}
 \mathcal L\varphi(\mu_*)
 \leq g_*\cdot\mathcal L\alpha(\mu_*)-1.
\end{equation}
\end{theorem}

Here ``touches from below'' means equality at the indicated point and the
inequality \(\varphi\leq F_\beta\circ\alpha\) in a neighborhood.  No
derivatives of \(F_\beta\) are assumed.  A supporting supergradient is a
vector such that
\[
F_\beta(q)\leq F_\beta(\alpha(\mu_*))+g_*\cdot(q-\alpha(\mu_*))
 \quad\hbox{for every source law }q.
\]
Such vectors exist at positive laws by Proposition~\ref{prop:values}.

Write \(h=F_\beta\), \(\mu_0=\mu_*\), \(\alpha_0=\alpha(\mu_0)\), and
\[
 \ell(q)=h(\alpha_0)+g_*\cdot(q-\alpha_0).
\]
For a finite construction with law \(q\) and utility \(u\), its
\emph{supporting-plane cost} is \(\ell(q)-u\).  It is nonnegative because
\(u\leq h(q)\leq\ell(q)\).  This cost measures lost utility; it is not
an additional auxiliary rank charge.

\subsection{Locally readable target windows}

We require many target indices with no nonconstant three-term arithmetic
progression.  We use the integer form of the progression-free set
construction of Behrend \cite{behrend}, proved in
Lemma~\ref{val:ap-set}.  Using such labels to separate tensor products
follows the progression-free extraction method of
Coppersmith--Winograd \cite[Sections~4--7]{cw}; the count-window
restriction needed here is given explicitly below.

\begin{lemma}\label{ap:behrend}
For all sufficiently large integers \(L\), there is a set
\(J\subseteq\{1,\ldots,L\}\) with no nonconstant three-term arithmetic
progression and
\[
 \log|J|\geq\log L-O(\sqrt{\log L}).
\]
\end{lemma}
\begin{proof}
This is the integer assertion of Lemma~\ref{val:ap-set}.
\end{proof}

For a length-\(n\) word with grades \(g_1,\ldots,g_n\), its centered
grade sum is \(\sum_{t=1}^n(g_t-\mu_0)\).  Use the targets
\begin{equation}\label{ap:target}
 y_j^{\mathrm{grade}}=(j,j,-2j)\sqrt n/L,
 \qquad j\in J,
\end{equation}
and coordinate windows with radius
\(\rho=\sqrt n/(10L)\).  Each side reads its own centered sum.
Because the three centered sums add to zero, a retained triple whose
three local window labels are \(j_0,j_1,j_2\) satisfies
\[
 |j_0+j_1-2j_2|\leq3/10.
\]
The left side is an integer, and hence is zero.  The progression-free
property gives \(j_0=j_1=j_2\).  Windows are disjoint on every side,
including the side whose target spacing is twice as large.  Thus a
nonzero retained tensor has a shared direct target label.

In a physical orientation bank \(\sigma\in\mathfrak S_3\), the mean
\(\mu_0\), the target vector, and the local windows are all permuted
by \(\sigma\).  The coordinate-sum test is expressed in canonical
coordinates.  A physical side reads the corresponding
permuted coordinate.  The coordinate-sum argument is unchanged, so
the target remains shared.  There is no averaging of the canonical
grade covariance across orientations in this operation.

\subsection{Exact finite compilation}

At each target and prefix, a control specifies a probability law on the
source alphabet and a cost.  A
\emph{permitted path} chooses at every step a letter in the prescribed
law's support.  In the next lemma, every node law and utility must be
supplied by an earlier finite strong realization.  The lemma turns these
realizations into tensor maps when every permitted path ends inside its
assigned target windows.
It is useful to state it for either alphabet, and with a general coefficient
\(\lambda>0\) on the direct-label logarithm.  For the six-state score
\(D-A+(\beta/3)M\) this coefficient is \(\lambda=1\); for the
W score \(\gamma(D-A)+M\) it is \(\lambda=\gamma\).

The exact populations and multinomial entropy estimates use the method
of types \cite[Section~II]{csiszar1998}.  The proof also establishes
the local tensor maps, their uniform auxiliary supply, and nonvanishing
of every counted output.

\begin{lemma}[Compiling a finite tree]\label{lem:compile}
Fix a finite horizon \(n\), a finite set of disjoint shared terminal
windows \(J\), and a rational target law \(a\).  At every target and
prefix node \(v=(j,h)\), prescribe an earlier finite strong realization
with exact canonical law \(q_v\), utility \(u_v\), and uniform auxiliary
and byproduct dimensions.  Assume that every path permitted by these
laws belongs to its target windows.  Put
\[
 w_{j,h}=a_j\prod_{t=1}^{|h|}q_{j,h_{<t}}(h_t),
 \qquad
 \overline q=\frac1n\sum_{|h|<n,\ j\in J}w_{j,h}q_{j,h}.
\]
Then finite strong realizations with exact law \(\overline q\) attain,
in the limit of their outer bank sizes,
\begin{equation}\label{ap:compile-score}
 u_{\mathrm{out}}
 \geq\frac{\lambda H(a)}n
      +\frac1n\sum_{|h|<n,\ j\in J}w_{j,h}u_{j,h}.
\end{equation}
The limiting notation only removes a multinomial type-count loss;
every member of the family is a finite tensor construction using one
fixed auxiliary tensor and nonzero counted outputs.
\end{lemma}

\begin{proof}
Use \(B\) length-\(n\) macroblocks in each of the six physical
orientations.  Thus the total number of source occurrences is \(6Bn\).
First apply the local terminal-window projections.  On any surviving
term, the preceding argument makes the target of each macroblock a
shared label.  On every orientation restrict the target-label word to
the exact type \(a\).  The number of such words is
\[
 \binom B{(Ba_j)_{j\in J}},
\]
with logarithm \(BH(a)+O(\log B)\), since \(J\) is fixed.  The
six orientations have independently chosen target words of this type.

For each shared target assignment, process sites in order.  At the node
\(v=(j,h)\), there are exactly
\[
 b_v=Bw_v
\]
macroblocks in every orientation having that target and canonical prefix.
Apply its prescribed strong realization to their next occurrences, as an
equal collection across the six orientations.  If that realization consumes
\(N_v\) occurrences in each orientation, tile \(b_v/N_v\) copies of
it.  Its exact per-bank histogram guarantees that the number of these
macroblocks receiving next canonical letter \(i\) is
\begin{equation}\label{ap:prefix-count}
 b_{j,hi}=b_{j,h}q_{j,h}(i).
\end{equation}
These numbers do not depend on which of its direct output labels occurs.
That label does identify, on every side, exactly which macroblocks enter
each child node, so the next maps may address the corresponding original
occurrence positions conditionally on the shared direct label.

All the numbers \(w_v\) and \(q_v\) are rational.  There are finitely
many nodes and finitely many integers \(N_v\).  Taking \(B\) through
common multiples of their denominators and the necessary \(N_v\)'s
makes every nonzero \(b_v\) an exact tiling.  Zero populations require
no call.  This proves the induction \eqref{ap:prefix-count} throughout
the tree.  It also shows that every branch uses the same number of every
raw auxiliary tensor.  Append their product once.  Apply the appropriate
maps inside each already direct branch; there is no multiplication of
auxiliary rank by the number of branch labels.  Uniformity of the raw
byproduct dimensions makes the final byproduct dimensions uniform as
well.  The byproducts remain separate tensor factors.

We justify carefully why the earlier count restrictions do not invalidate
the raw calls.  Definition~\ref{def:strong} and the pointwise construction
in Lemma~\ref{alg:payload} give a coordinatewise
identity for every correlated unresolved payload, preserving the original
source coordinates.  Conditional maps can address the listed original
positions without changing the source word.  Their identities therefore
commute with the previously applied local count projectors.  The same
statement holds for a degeneration after taking its leading coefficient.
Equivalently, compute the unfiltered raw identity first and then apply
the original count projections to its designated fibers.

This identity alone would allow a formal zero fiber, which must not be
counted.  Here every complete designated word has, in each macroblock,
a path permitted by the chosen controls.  It consequently satisfies that
macroblock's target windows.  Its local raw words satisfy the support
conditions of the corresponding earlier strong realizations.  Starting from
the independent source tensor, every such complete source word is present,
and it survives all of these predicates.  Thus every designated output
counted here is nonzero.  This also proves that every target assignment
of the retained exact type supplies all the claimed raw outputs, even
though the initial projection may correlate the unprocessed suffixes.

Bank labels are fixed physical orientations throughout.  At a node the
child realization supplies law \(q_v\) in canonical coordinates in each
bank; its physical law in bank \(\sigma\) is \(\sigma q_v\).
Both the prefix classification and the window centering use that same
conversion.  This AP compilation introduces no additional relative
permutation of a site's source orientation: all internal physical
permutations of a raw recipe are already included in its strong realization
identity and its exact per-bank type.  The population assigned to an
original input bank is not redefined after observing an output label.
Summing the node contributions therefore gives the exact final canonical
law \(\overline q\) in every bank.

For clarity about the probability notation, choose one macroblock uniformly
from any fixed final output branch, within a fixed canonical orientation.
The exact counts give
\[
 \Prob(j,h)=w_{j,h},
 \qquad
 \Prob(\text{next letter}=i\mid j,h)=q_{j,h}(i).
\]
This is a counting identity.  Different macroblocks need not be independent.
The law of one sampled block is exactly the law used in the tree's
expected-cost calculation, so its single-step canonical covariance is
the covariance of \(q_{j,h}\), not a symmetrization across banks.

It remains to count multiplicities and charges.  If the child realization
at \(v\) has direct multiplicity \(L_v\), rank bound \(R_v\), and
volume \(V_v\), its rates are normalized by \(6N_v\).  Its
\(b_v/N_v\) copies therefore contribute precisely \(w_v/n\) times
each raw rate to the rate per \(6Bn\) source occurrences.  The target
types contribute
\[
 \frac{6\log\binom B{(Ba_j)_j}}{6Bn}
  =\frac{H(a)}n+O\left(\frac{\log B}{B}\right)
\]
to the direct rate and nothing to the auxiliary rate.  This proves
\eqref{ap:compile-score} in the limit.

No extra transcript entropy has been counted: the target is recovered
uniquely from the final original word's disjoint count windows.  Once
the target and earlier prefixes are fixed, injectivity of each child
realization recovers its direct label from its designated original subword.
Induction therefore makes the complete output label inject into the
complete source word.  The coordinatewise identity, exact bank law,
uniform dimensions, and nonvanishing just proved are precisely the strong
realization invariants.  Only finite compositions and tensor products have
been used for every admissible finite \(B\).
\end{proof}

For any affine functional \(\ell\), define the node cost by
\(c_v=\ell(q_v)-u_v\).  The prefix weights sum to one at each time,
so the averaged law in the lemma satisfies
\[
 \frac1n\sum_v w_v\ell(q_v)=\ell(\overline q).
\]
Thus \eqref{ap:compile-score} has the equivalent form
\begin{equation}\label{ap:compile-gain}
 u_{\mathrm{out}}
 \geq\ell(\overline q)
       +\frac{\lambda H(a)-\mathbb E\sum_{t=1}^n c_t}{n}.
\end{equation}
The expectation denotes the finite sum \(\sum_v w_vc_v\).
For the six-state supporting plane fixed above, \(\lambda=1\) and
every node cost is nonnegative.  To contradict that plane, it is therefore
enough to send every permitted path into its window while keeping the
expected cost strictly below \(H(a)\).  The remaining argument constructs
such controls from a failed lower test.

\subsection{A lower test supplies inexpensive small drifts}
\label{ap:local-cost}

At a lower contact, \(\ell\circ\alpha-\varphi\) has a local minimum
zero.  Thus its first derivative vanishes and
\begin{equation}\label{ap:q0}
 Q_0=g_*\cdot D^2\alpha(\mu_0)-D^2\varphi(\mu_0)\succeq0.
\end{equation}
If \eqref{ap:lower-test} fails, \(\Tr(Q_0C(\alpha_0))<1\).
Choose \(\varepsilon>0\) so small that
\begin{equation}\label{ap:tau}
 Q=Q_0+\varepsilon I\succ0,
 \qquad \tau=\Tr(QC(\alpha_0))<1.
\end{equation}
Taylor's theorem and \(\varphi\leq h\circ\alpha\) give, after reducing
the neighborhood if necessary,
\begin{equation}\label{ap:local-cost-bound}
 0\leq c(v):=\ell(\alpha(\mu_0+v))-h(\alpha(\mu_0+v))
 \leq \frac12v^\top Qv.
\end{equation}
By the definition of the closed attainable set, the law and utility in
this display can be approximated simultaneously, to any prescribed positive
accuracy, by a finite earlier construction.  We postpone these finitely
many approximations until the entire control tree has been fixed.

Transform centered grades by \(Q^{1/2}\).  A grade \(g\) then gives the
increment \(X=Q^{1/2}(g-\mu_0)\).  A small prescribed mean \(m\) is
obtained by taking the law
\(\alpha(\mu_0+Q^{-1/2}m)\).  We call these chart laws soft controls;
their ideal cost satisfies
\begin{equation}\label{ap:drift-cost}
 c(m)\leq |m|^2/2.
\end{equation}
There is a fixed bound \(B>0\) on the magnitude of every increment, and
\begin{equation}\label{ap:noise-limit}
 \E|X|^2\longrightarrow\tau\quad\hbox{as }m\longrightarrow0.
\end{equation}
The limit concerns the second moment, which equals the covariance trace
at mean zero.

We also allow exact singleton controls.  The point \(\mu_0\) is in the
relative interior of the grade convex hull, so there is a \(\kappa>0\)
such that for every unit vector \(u\in\mathcal P\), one singleton
increment satisfies \(u\cdot X\leq-\kappa\).  Its utility is zero.
The costs of the finitely many singleton choices have a finite common
upper bound \(C_*\).

\subsection{A control that retains every path}

We now construct controls for the bounded increments just described whose
cost is smaller than the target-label gain in \eqref{ap:compile-gain}.
The terminal guarantee must hold for every permitted path.  Soft controls
will be made safe for every letter in the entire source alphabet, so later
rational approximation cannot create an unsafe new outcome.

\begin{lemma}[Shrinking-corridor control]\label{ap:corridor}
Assume that every \(|m|\leq m_0\) is available as a soft mean, with
cost at most \(|m|^2/2\), bounded increments \(|X|\leq B\), and the
limit \eqref{ap:noise-limit}.  Assume also the singleton controls above.
Fix \(c>0\) and a target bound \(C_y>0\).  For every \(\eta>0\)
there is a constant \(K_\eta\), independent of \(L\), with the following
property.  For each integer \(L\geq2\), all sufficiently large finite
horizons \(n\), and every target \(y\) with \(|y|\leq C_y\sqrt n\),
there is a finite control tree such that every permitted path satisfies
\[
 \left|\sum_{t=1}^n X_t-y\right|<w,
 \qquad w=c\sqrt n/L,
\]
and its expected total cost is at most
\begin{equation}\label{ap:total-cost}
 \E\sum_{t=1}^n c_t\leq(\tau+2\eta)\log L+K_\eta.
\end{equation}
The expectation is only an accounting device for the tree's transition
laws; the terminal inclusion holds for every path.
\end{lemma}

\begin{proof}
Reduce \(m_0\) if necessary so that all its means lie in the available
local chart.  Choose a wall cutoff \(d_0\) such that
\[
 d_0\geq16B,
 \qquad d_0m_0\geq1024B^2.
\]
Choose \(a\) sufficiently large that
\[
 a\geq1024B^2,
 \qquad
 a\log\left(1+\frac{\kappa}{8d_0}\right)\geq2C_*+2.
\]
Finally choose \(D_0\) so large that
\begin{equation}\label{ap:constants}
 D_0>2C_y+1,
 \qquad D_0\geq B,
 \qquad \frac{aB^2}{D_0^2}<\eta/4.
\end{equation}
The numerical factors are convenient slack, not optimized constants.
All these choices are independent of \(L\) and \(n\).

If \(r\) steps remain, put
\[
 s=r+w^2/D_0^2,
 \qquad R_s=D_0\sqrt s,
 \qquad
 J_s(e)=\frac{|e|^2}{2s}
       -a\log\left(1-\frac{|e|^2}{D_0^2s}\right),
\]
where \(e\) is the current displacement from the target.  The permitted
region is \(|e|<R_s\).  Initially \(e=-y\), so this region contains
the starting point with a fixed relative margin.  At the terminal time
its radius is \(w\).

Write \(\delta=R_s-|e|\).  If \(\delta\geq d_0\), choose the soft
mean
\begin{equation}\label{ap:feedback}
 m=-\operatorname{clip}_{m_0}
   \left(\frac e s+\frac{2ae}{R_s^2-|e|^2}\right),
\end{equation}
where clipping truncates the vector's length to at most \(m_0\).
If \(\delta<d_0\), choose an exact singleton with inward radial
component at least \(\kappa\).

First verify support.  In the soft region every increment in the entire
alphabet has magnitude at most \(B<d_0\), while
\(R_s-R_{s-1}=O(s^{-1/2})\).  Thus every outcome stays in the next
region once the minimum \(s\) is large.  In the singleton region,
\(|e|\to\infty\) uniformly as that minimum grows, and the inward
radial component gives
\[
 |e+X|\leq |e|-\kappa/2
\]
for all sufficiently large \(s\).  It too stays in the next region,
whose shrinkage is eventually at most \(\kappa/4\).  Induction proves
the support assertion at every time.  At fixed \(L\), the minimum
\(s=w^2/D_0^2=c^2n/(D_0^2L^2)\) tends to infinity with \(n\), so
these statements are uniform over the finite horizon.

We next prove the uniform one-step estimate
\begin{equation}\label{ap:one-step}
 c_{\rm step}+\E J_{s-1}(e+X)-J_s(e)
 \leq\frac{\tau/2+\eta}{s}.
\end{equation}
Here \(c_{\rm step}\) is the cost of the chosen control at the
current state; the constant \(c\) in \(w=c\sqrt n/L\) remains fixed.
There are three regimes.  We give their estimates explicitly, and then
explain uniformity.

\paragraph{Interior regime.}
Suppose along a sequence with \(s\to\infty\) that
\(t=|e|^2/(D_0^2s)\) stays bounded away from one.  The drift in
\eqref{ap:feedback} is then unclipped and tends to zero.  Taylor
expansion in space and in the time variable, with remainder \(o(s^{-1})\),
is valid uniformly on any such smaller region.  Indeed the third spatial
derivatives are \(O(s^{-3/2})\), and the needed time and mixed
derivatives have their corresponding smaller orders, while \(|X|\leq B\).
Using \(c_{\rm step}\leq |m|^2/2\) and \(m=-\nabla J_s\), the leading temporal
and drift contributions, after multiplication by \(s\), sum to
\[
 -\frac{at}{1-t}
 -\frac{2a^2t}{D_0^2(1-t)^2}.
\]
For completeness, the temporal contribution is
\(D_0^2t/2+at/(1-t)\); the combined drift and cost contribution is
\(-D_0^2t/2-2at/(1-t)-2a^2t/[D_0^2(1-t)^2]\).
The quadratic noise gives \(\tau/2+o(1)\) by
\eqref{ap:noise-limit}.  The remaining barrier noise is at most
\[
 \frac{aB^2}{D_0^2(1-t)}
 +\frac{2aB^2t}{D_0^2(1-t)^2}.
\]
The second of these terms is absorbed by the negative term quadratic
in \(a\), since \(a\geq B^2\).  The first, together with
\(-at/(1-t)\), is at most \(aB^2/D_0^2\), since \(D_0^2\geq B^2\).
Consequently the limsup of the left side of \eqref{ap:one-step}
multiplied by \(s\) is at most \(\tau/2+\eta/4\).

\paragraph{Soft control near the wall.}
Now suppose \(\delta\geq d_0\) and \(\delta/R_s\to0\).  The drift
is radial and inward, and
\begin{equation}\label{ap:wall-mean}
 |m|=(1+o(1))\min\{m_0,a/\delta\}.
\end{equation}
This follows because the radial length before clipping is
\(|e|/s+2a|e|/[\delta(R_s+|e|)]\), whose second term is
\((1+o(1))a/\delta\) and whose first is smaller by
\(O(\delta/\sqrt s)\).  Define
\[
 z_*=\frac{D_0^2+2e\cdot X+|X|^2}{R_s^2-|e|^2}.
\]
The logarithmic part of the potential changes by
\[
 a\{-\log(1-z_*)+\log(1-1/s)\}.
\]
For large \(s\), \(|z_*|\leq3B/\delta\leq1/2\).  The inequality
\(-\log(1-z)\leq z+2z^2\) for \(|z|\leq1/2\) therefore gives an
upper bound
\begin{equation}\label{ap:wall-barrier}
 -(1-o(1))\frac{a|m|}{\delta}
 +\frac{20aB^2}{\delta^2}
 +o\left(\frac{a|m|}{\delta}\right).
\end{equation}
To see the orders here, the mean of the linear part is
\[
 \E z_*
 =-\frac{2|e|}{R_s+|e|}\frac{|m|}{\delta}
   +O\left(\frac{D_0^2+B^2}{R_s\delta}\right),
\]
and \(\E z_*^2\leq9B^2/\delta^2\) for large \(s\).
The displayed error in the mean is
\(o(|m|/\delta)\) by \eqref{ap:wall-mean} and
\(\delta/R_s\to0\).  We used the harmless constant 20 to absorb
the second-moment bound.

The cost is at most
\((1+o(1))a|m|/(2\delta)\).  The quadratic part of the potential
has increment
\[
 \frac{|e|^2}{2s(s-1)}+\frac{e\cdot m}{s-1}
                       +\frac{\E|X|^2}{2(s-1)}.
\]
Its middle term is nonpositive, and the other terms are \(O(s^{-1})\),
which is \(o(a|m|/\delta)\) in this regime.  Finally,
\[
 \delta|m|\geq(1-o(1))\min\{d_0m_0,a\}
                  \geq(1-o(1))1024B^2.
\]
Thus the noise term in \eqref{ap:wall-barrier} is strictly smaller
than the remaining negative drift term.  The left side of
\eqref{ap:one-step} is negative for all sufficiently large \(s\)
along any sequence in this regime.

\paragraph{Singleton control near the wall.}
If \(\delta<d_0\), the support estimate already proved gives the new
wall distance \(\delta'\geq\delta+\kappa/4\).  Using
\(R_s^2-|e|^2=\delta(R_s+|e|)\), and noting that the other radius
factors tend to one uniformly, the logarithmic potential decreases by
at least
\[
 a\log\left(1+\frac{\kappa}{8d_0}\right)
\]
for sufficiently large \(s\).  This pays the cost \(c_{\rm step}\leq C_*\)
with fixed slack by the choice of \(a\).  The quadratic-potential
increment is \(O(s^{-1/2})\) and hence cannot consume that slack.
The left side of \eqref{ap:one-step} is again negative.

These estimates imply a uniform threshold for \eqref{ap:one-step}.
Otherwise a sequence of violating states with \(s\to\infty\) would
have a subsequence on which \(t\) stays away from one, or tends to one.
The former contradicts the interior estimate with its strict
\(3\eta/4\) margin.  The latter has a further subsequence in one of
the two wall regimes, where the left side is negative.  This also proves
that the threshold works for every time in a sufficiently large horizon.

The initial value of \(J_s(-y)\) is bounded by a constant depending
only on the chosen control constants and \(C_y\), not on \(L\).
The final value is nonnegative.  Summing \eqref{ap:one-step} and
telescoping therefore bounds the expected cost by
\[
 (\tau/2+\eta)\sum_{r=1}^n
   \frac1{r+w^2/D_0^2}+O_\eta(1)
 \leq(\tau/2+\eta)
   \log\left(1+\frac{D_0^2L^2}{c^2}\right)+O_\eta(1).
\]
Since the logarithm is at most \(2\log L+O(1)\), this is
\eqref{ap:total-cost}.  The factor two in this logarithmic time ratio
cancels the factor \(1/2\) in the one-step covariance contribution.
\end{proof}

\subsection{The contradiction and the order of limits}

\begin{proof}[Proof of Theorem~\ref{thm:ap}]
Suppose the conclusion fails and choose \(Q\) and \(\tau<1\) as in
\eqref{ap:tau}.  Choose \(\eta>0\) with \(\tau+2\eta<1\).
Transform the grade-sum targets in \eqref{ap:target} by \(Q^{1/2}\).
Their norms are at most \(C_y\sqrt n\) for a fixed \(C_y\), independent
of \(L\) and \(n\).  Choose \(c>0\) so that
\(c\|Q^{-1/2}\|\leq1/10\).  A transformed ball of radius
\(c\sqrt n/L\) then lies inside all three original coordinate windows.
Fix the resulting constants of Lemma~\ref{ap:corridor}.
For each \(L\), Lemma~\ref{ap:behrend}
gives a target set \(J\).  Give its targets equal probabilities, so
\(a_j=1/|J|\) and \(H(a)=\log|J|\).  The average ideal steering
cost is bounded by
\[
 \E\sum c_t\leq(\tau+2\eta)\log L+K_\eta.
\]
Choose \(L\) so large that
\begin{equation}\label{ap:strict-gain}
 G:=\log|J|-\{(\tau+2\eta)\log L+K_\eta\}>0,
\end{equation}
using the uniform upper bound on the expectation.  Then choose a
sufficiently large finite \(n\) for all the control estimates and fix
the complete finite control tree supplied by Lemma~\ref{ap:corridor}
for every target.

For each soft node the law is the appropriate chart law and the ideal
utility is \(h\) at that law.  Exact singleton nodes stay exact.
Include all source letters as possible children of every soft node,
whether or not an ideal law assigns them positive probability, and give
every such child the safe continuation supplied by the controller.
Thus the pathwise window property is independent of small perturbations
of the soft transition probabilities.

There are now only finitely many soft nodes.  Approximate each one's
ideal law and utility by an earlier finite strong realization from the
closed attainable class.  The finite realization has a rational exact
law, which need not lie exactly on the smooth chart.  This causes no
difficulty.  Its actual supporting-plane cost is
\(\ell(q_v)-u_v\geq0\), and approaches the ideal cost as the law and
utility errors tend to zero.  In a finite tree, all prefix probabilities,
the expected total cost, and the averaged source law depend continuously
on these finitely many data.  Choose the approximation errors so small
that the actual expectation is at most the ideal expectation plus
\(G/2\).  No accuracy bound uniform in \(n\) is needed, since this
finite \(n\) has already been fixed.

Apply Lemma~\ref{lem:compile} to these finite realizations.  Let
\(\overline\alpha\) be its exact averaged law.  The supporting affine
terms sum linearly, because at every time the prefix weights sum to one.
Thus its limiting utility satisfies
\begin{align*}
 u_{\mathrm{out}}
 &\geq\frac{\log|J|}n
   +\frac1n\sum_v w_v\{\ell(q_v)-c_v\}\\
 &=\ell(\overline\alpha)
   +\frac{\log|J|-\E\sum_t c_t}{n}
 \geq\ell(\overline\alpha)+\frac{G}{2n}.
\end{align*}
Taking the finite outer bank size sufficiently large makes its remaining
type-count loss less than \(G/(4n)\).  We have therefore constructed
an actual finite admissible realization whose utility exceeds
\(\ell(\overline\alpha)\), contradicting the supporting-plane bound
\(h\leq\ell\).

The order of choices was: local control constants; a finite target set;
a sufficiently large finite horizon and its full tree; finitely many
finite approximating raw recipes; and finally a sufficiently large outer
common multiple.  Each approximant has finite inductive depth, so their
maximum depth is finite and the new macro uses only earlier constructions.
The potentially very large common multiple affects existence neither
of the maps nor of the strict excess.  No infinite-depth realization has
been applied to a finite tensor.  The contradiction proves
\(\Tr(Q_0C)\geq1\), which is exactly \eqref{ap:lower-test}.
\end{proof}

Write \(U_\gamma=U^W_\gamma\) for the three-state attainable value.

\begin{corollary}[Lower tests for the W value]\label{cor:w-ap}
Fix \(\gamma\geq1\).  On the relative interior of the three-letter
simplex, put
\[
 C_W(q)=\diag(q)-qq^\top,
 \qquad \mathcal L_W=C_W(q):D^2.
\]
Every \(C^2\) lower test \(\psi\) for \(U_\gamma\) satisfies
\begin{equation}\label{ap:w-test}
 \mathcal L_W\psi\leq-\gamma
\end{equation}
at its contact point.
\end{corollary}

\begin{proof}
Use W grades \(e_0,e_1,e_2\), whose means equal their law \(q\).
The chart is affine, its covariance is \(C_W(q)\), and the concavity,
supporting planes, and boundedness needed above hold by
Proposition~\ref{prop:values}.  The argument from a failed lower test
again gives a quadratic drift cost and a trace \(\tau\).  All control
and target-window proofs apply because the centered grades sum to zero
and span the same plane.  The sole change in the utility accounting is
that a direct target contributes \(\gamma\log|J|\).  Lemma~\ref{lem:compile}
uses \(\lambda=\gamma\), while the hard-wall cost is still
\((\tau+2\eta)\log L+O_\eta(1)\).  Thus a contradiction occurs if
\(\tau<\gamma\), proving \eqref{ap:w-test}.  No pre-existing payload
volume enters this utility.
\end{proof}

The following interior-maximum argument is the strict-comparison
strategy of viscosity-solution theory
\cite[Section~5.C]{crandallishiilions1992}.  Its short proof uses the
lower-test inequalities just established.

\begin{corollary}[Strict comparison]\label{ap:comparison}
Let \(\Omega\) be a bounded domain in the mean plane with compact
closure, and let a chart \(\alpha\) as in Theorem~\ref{thm:ap} extend
continuously to \(\overline\Omega\).  Suppose
\(f\in C(\overline\Omega)\cap C^2(\Omega)\),
\(f\leq F_\beta\circ\alpha\) on \(\partial\Omega\), and, for every
\(\mu\in\Omega\) and every supporting supergradient \(g\) there,
\[
 \mathcal L f(\mu)>g\cdot\mathcal L\alpha(\mu)-1.
\]
Then \(f\leq F_\beta\circ\alpha\) throughout \(\overline\Omega\).
For the affine W chart, the analogous conclusion holds for \(U_\gamma\)
when \(\mathcal L_W f>-\gamma\).
\end{corollary}

\begin{proof}
The value functions are continuous by Proposition~\ref{prop:values}.
If \(f-F_\beta\circ\alpha\) has a positive maximum \(c\), compactness
and the boundary inequality place a maximizing point in the interior.
The function \(f-c\) touches the value from below there.  Its derivatives
equal those of \(f\), contradicting Theorem~\ref{thm:ap}.  The W proof
uses Corollary~\ref{cor:w-ap} in the same way.
\end{proof}

\subsection{Controlling the supporting-plane term}

The lower-test inequality retains \(g\cdot\mathcal L\alpha\) when the
chart is not exactly harmonic.  Although the supporting supergradient is
unknown, concavity and a uniform value bound control this term.  The
certificate will apply the following estimate where every chart coordinate
has a positive lower bound.

\begin{lemma}[Supporting-plane chart error]\label{ap:chart-error}
Suppose \(0\leq h\leq K\) is a concave value on the full source
simplex, \(\alpha\) is an interior law, and \(g\) is a supporting
supergradient there.  For a probability-law chart through \(\alpha\),
\[
 |g\cdot\mathcal L\alpha|
 \leq K\max_i\frac{|\mathcal L\alpha_i|}{\alpha_i}.
\]
In particular, if \(\alpha_i\geq m>0\) and
\(\|\mathcal L\alpha\|_\infty\leq\varepsilon\), the right side
is at most \(K\varepsilon/m\).
\end{lemma}

\begin{proof}
Add a multiple of the all-ones vector to \(g\) so that
\(g\cdot\alpha=h(\alpha)\).  This leaves its supergradient inequality
on the simplex unchanged, and does not change \(g\cdot\mathcal L\alpha\)
because \(\sum_i\alpha_i=1\).  The resulting supporting plane is
\(q\mapsto g\cdot q\).  Evaluating it at simplex vertices and using
\(h\geq0\) gives \(g_i\geq0\).  Therefore
\[
 |g\cdot\mathcal L\alpha|
 \leq\sum_i g_i|\mathcal L\alpha_i|
 \leq\left(\max_i\frac{|\mathcal L\alpha_i|}{\alpha_i}\right)
       \sum_i\alpha_i g_i
 \leq K\max_i\frac{|\mathcal L\alpha_i|}{\alpha_i}.
\]
\end{proof}
